\begin{pastexam}[label={px:gw3-097-ds-07}]{關務三等資料處理 97 年 第七題}
（申論；15 分）

例舉說明何謂 2 元樹、AVL 樹、2-3-4 樹，並說明其異同點。（15 分）
\end{pastexam}

\begin{pastsoln}
下面三個例子都用同一組鍵，依序插入 $10, 20, 30, 40, 50, 60, 70$，方便比較。

\textbf{二元樹（2 元樹）}：每個節點最多有兩個子節點，分成左子、右子（有順序之分）的樹；空樹也算
二元樹（資結書第~7.1~節）。 二元樹本身對鍵的大小沒有規定；拿來搜尋時用的是\textbf{二元搜尋樹}：
左子樹的鍵都比節點小、右子樹的都比節點大。 普通的二元搜尋樹不做任何調整，依序插入遞增的鍵會變成
一條往右的鏈，高度 $n - 1$，搜尋退化成 $O(n)$（下圖左）。

\textbf{AVL 樹}：一種二元搜尋樹，額外要求\textbf{每個節點左右子樹的高度差不超過 $1$}。 插入或刪除後若
有節點高度差變成 $2$，用旋轉（LL、RR、LR、RL）恢復，因此高度一定是 $O(\lg n)$（約 $1.44 \lg n$
以內），搜尋、插入、刪除都是 $O(\lg n)$（資結書第~10.2~節）。 同一組鍵插入時，$30$、$50$、$60$、$70$
各引起一次 RR 左旋，最後是一棵完全平衡的樹（下圖中）。

\textbf{2-3-4 樹}：一種多路搜尋樹，每個節點存 $1$、$2$ 或 $3$ 個鍵，分別有 $2$、$3$、$4$ 個子節點
（2-節點、3-節點、4-節點）；節點內的鍵由小到大，第 $i$ 棵子樹的鍵介於第 $i-1$ 和第 $i$ 個鍵之間；
而且\textbf{所有葉節點都在同一層}。 插入時遇到滿的 4-節點就把中間的鍵送上父節點、分裂成兩個
2-節點，樹只會從根往上長高，所以永遠平衡，高度介於 $\log_4(n+1)$ 和 $\lg(n+1)$ 之間（資結書
第~11.1.5~節）。 同一組鍵由上而下插入的結果見下圖右。
\begin{center}
\begin{tikzpicture}[
  v/.style={circle, draw, minimum size=6.5mm, inner sep=0pt, font=\small},
  r/.style={draw, minimum height=5.5mm, inner sep=2pt, font=\small}]
\begin{scope}[x=0.45cm, y=0.62cm]
\node[font=\small] at (3,1.3) {二元搜尋樹};
\foreach \k [count=\i from 0] in {10,20,30,40,50,60,70} \node[v] (b\i) at (\i,-\i*0.62) {\k};
\foreach \i [evaluate=\i as \j using int(\i+1)] in {0,...,5} \draw (b\i)--(b\j);
\end{scope}
\begin{scope}[xshift=5.6cm, level distance=10mm,
  level 1/.style={sibling distance=22mm}, level 2/.style={sibling distance=11mm}]
\node[font=\small] at (0,0.8) {AVL 樹};
\node[v] {40}
  child {node[v] {20} child {node[v] {10}} child {node[v] {30}}}
  child {node[v] {60} child {node[v] {50}} child {node[v] {70}}};
\end{scope}
\begin{scope}[xshift=10.8cm, level distance=10mm, level 1/.style={sibling distance=13mm}]
\node[font=\small] at (0,0.8) {2-3-4 樹};
\node[r] {20 | 40}
  child {node[r] {10}} child {node[r] {30}} child {node[r] {50 | 60 | 70}};
\end{scope}
\end{tikzpicture}
\end{center}

\textbf{相同點。}
\begin{itemize}[nosep]
\item 三者都是樹狀結構，用指標連接節點；二元搜尋樹、AVL 樹、2-3-4 樹都滿足「左小右大」的搜尋樹
性質，中序走訪會由小到大列出所有鍵，搜尋都是從根往下、每層比較後選一條路走。
\item AVL 樹和 2-3-4 樹都是\textbf{平衡搜尋樹}，保證高度 $O(\lg n)$，搜尋、插入、刪除最壞都是
$O(\lg n)$。
\end{itemize}

\textbf{相異點。}
\begin{center}
\small
\begin{tabular}{p{2.3cm}p{3.4cm}p{3.6cm}p{3.6cm}}
\toprule
 & 二元（搜尋）樹 & AVL 樹 & 2-3-4 樹 \\
\midrule
每節點的鍵／子節點 & 1 個鍵、最多 2 個子 & 1 個鍵、最多 2 個子 & 1--3 個鍵、2--4 個子（葉除外） \\
平衡條件 & 無 & 每個節點左右高度差 $\le 1$ & 所有葉節點同一層 \\
最壞高度 & $n - 1$ & 約 $1.44 \lg n$ & $\lg(n+1)$ \\
維持平衡的方法 & 不維持 & 旋轉 & 節點分裂、合併、向兄弟借鍵 \\
最壞搜尋時間 & $O(n)$ & $O(\lg n)$ & $O(\lg n)$ \\
\bottomrule
\end{tabular}
\end{center}
另外，2-3-4 樹的每個節點可以用一個黑節點加上 $0$ 到 $2$ 個紅節點表示，變成一棵紅黑樹；所以 2-3-4 樹
也可以看成紅黑樹的「多路」寫法（資結書第~10.4~節）。
\end{pastsoln}
